% =====================================================================
\chapter{Convexidad}
\label{tema:5}
% =====================================================================

En una función convexa, todo mínimo local es global, de modo que un método de descenso no puede quedar atrapado en un mínimo malo. Este tema define los conjuntos y las funciones convexas, demuestra las caracterizaciones de primer y segundo orden, establece las operaciones que conservan la convexidad y las aplica a la entropía cruzada. Termina explicando por qué la pérdida de una red profunda no es convexa. Se sigue~\cite[§3.1, §3.2 y §4.2]{bv} y las transparencias correspondientes de~\cite{ee364a}.

\section{Conjuntos y funciones convexas}

\begin{definicion}[Conjunto convexo]
Un conjunto $C\subseteq\R^{n}$ es \emph{convexo} si, para todo $a,b\in C$ y todo $\lambda\in[0,1]$, el punto $\lambda a+(1-\lambda)b$ pertenece a $C$: contiene el segmento que une dos cualesquiera de sus puntos.
\end{definicion}

\begin{definicion}[Función convexa]\label{def:convexa}
Sea $C\subseteq\R^{n}$ convexo. Una función $f\colon C\to\R$ es \emph{convexa} si
\[
f\bigl(\lambda a+(1-\lambda)b\bigr)\le\lambda f(a)+(1-\lambda)f(b)\qquad\text{para todo } a,b\in C,\ \lambda\in[0,1],
\]
y \emph{estrictamente convexa} si la desigualdad es estricta cuando $a\ne b$ y $\lambda\in(0,1)$. Se dice que $f$ es \emph{cóncava} si $-f$ es convexa.
\end{definicion}

Geométricamente, la cuerda que une $(a,f(a))$ con $(b,f(b))$ queda por encima de la gráfica (\cref{fig:convexa}). Son convexas, por ejemplo, las funciones afines $a\T x+b$ (con igualdad en la definición), toda norma (por la desigualdad triangular y la homogeneidad) y $x\mapsto e^{x}$ en $\R$.

\addtocounter{figure}{1}%

\begin{lema}[Restricción a rectas]\label{lem:rectas}
Sea $C$ abierto y convexo. Una función $f\colon C\to\R$ es convexa si y solo si, para todo $x\in C$ y todo $v\in\R^{n}$, la función de una variable $\varphi(t)=f(x+tv)$ es convexa en el intervalo $\{t:x+tv\in C\}$.
\end{lema}

\section{Condiciones de primer y segundo orden}

\begin{teorema}[Condición de primer orden]\label{teo:primer-orden}
Sea $C$ abierto y convexo y $f\colon C\to\R$ diferenciable. Entonces $f$ es convexa si y solo si
\[
f(b)\ge f(a)+\nabla f(a)\T(b-a)\qquad\text{para todo } a,b\in C.
\]
\end{teorema}

La condición de primer orden dice que el plano tangente en cualquier punto es una cota inferior global de $f$: la información local del gradiente proporciona información global.

\begin{teorema}[Condición de segundo orden]\label{teo:segundo-orden}
Sea $C$ abierto y convexo y $f\in C^{2}(C)$. Entonces $f$ es convexa si y solo si $H_f(x)\succeq 0$ para todo $x\in C$. Si $H_f(x)\succ 0$ para todo $x\in C$, entonces $f$ es estrictamente convexa.
\end{teorema}

\addtocounter{definicion}{1}%

\begin{observacion}
El recíproco de la segunda parte del \cref{teo:segundo-orden} es falso: $f(x)=x^{4}$ es estrictamente convexa en $\R$, pero $f''(0)=0$.
\end{observacion}

\section{Mínimos de funciones convexas}

\begin{teorema}[Mínimos locales y globales]\label{teo:local-global}
Sea $f\colon C\to\R$ convexa. Todo mínimo local de $f$ es un mínimo global. Si además $f$ es diferenciable y $C$ es abierto, un punto $x^{*}$ es mínimo global si y solo si $\nabla f(x^{*})=0$.
\end{teorema}

Para funciones no convexas, $\nabla f=0$ solo es una condición necesaria. Para funciones convexas es también suficiente, y además global. Por eso el descenso de gradiente con una tasa adecuada resuelve los problemas convexos, como la regresión lineal y la regresión logística o softmax del módulo~8.

\section{Operaciones que conservan la convexidad}

\begin{proposicion}\label{prop:operaciones}
\begin{enumerate}
  \item Si $f_1,\dots,f_k$ son convexas y $w_1,\dots,w_k\ge 0$, entonces $\sum_i w_if_i$ es convexa.
  \item Si $f\colon\R^{m}\to\R$ es convexa, $A\in\R^{m\times n}$ y $b\in\R^{m}$, entonces $g(x)=f(Ax+b)$ es convexa.
  \item Si $f_1$ y $f_2$ son convexas, $\max(f_1,f_2)$ es convexa.
\end{enumerate}
\end{proposicion}

\begin{proposicion}[Convexidad de log-suma-exp]\label{prop:lse-convexa}
La función $\LSE(z)=\ln\sum_{k=1}^{K}e^{z_k}$ es convexa en $\R^{K}$. Su hessiano es $H=\diag(p)-pp\T$ con $p=\softmax(z)$, y para todo $v\in\R^{K}$
\[
v\T Hv=\sum_i p_iv_i^{2}-\Bigl(\sum_i p_iv_i\Bigr)^{2}=\Var_{i\sim p}[v_i]\ge 0,
\]
con igualdad si y solo si $v$ es proporcional a $\uno$.
\end{proposicion}

La demostración de~\cite[§3.1]{bv} obtiene la misma desigualdad mediante la de Cauchy--Schwarz.

\begin{corolario}[Convexidad de la entropía cruzada]\label{cor:ce-convexa}
La entropía cruzada $L(z)=-z_c+\LSE(z)$ es convexa en los logits $z$. Para un vector $x$ fijo, $W\mapsto -\ln\softmax(Wx)_c$ es convexa en $W$, y la pérdida media de la regresión softmax sobre un conjunto de datos es convexa en sus parámetros.
\end{corolario}

Para $z=(2,1,0)\T$, el hessiano de la entropía cruzada tiene autovalores $0$, $0{,}119$ y $0{,}371$. El autovalor nulo corresponde a la dirección $\uno$, a lo largo de la cual la pérdida es constante.

\section{Las redes profundas no son convexas}

\begin{proposicion}[Simetría y convexidad]\label{prop:simetria}
Sea $L\colon\R^{N}\to\R$ convexa. El conjunto de sus mínimos globales es convexo. En consecuencia, si $\theta_1$ y $\theta_2$ son mínimos globales y $\tfrac12(\theta_1+\theta_2)$ no lo es, $L$ no es convexa.
\end{proposicion}

En una red con una capa oculta, intercambiar dos neuronas ocultas (junto con sus pesos de entrada y de salida) no cambia la función calculada. Si $\theta_1$ es un mínimo, la permutación $\theta_2$ también lo es, pero el punto medio promedia las dos neuronas y, en general, calcula otra función con pérdida mayor. Por la \cref{prop:simetria}, la pérdida no es convexa.

\addtocounter{definicion}{1}%

\addtocounter{figure}{1}%

\addtocounter{definicion}{1}%

